\documentclass{article}
\usepackage{amsmath,amssymb,amsthm,algorithm,algorithmic}
\usepackage{fullpage}
\newcommand{\E}{\mathbb{E}}
\newcommand{\R}{\mathbb{R}}
\newcommand{\norm}[1]{\left\|#1\right\|}
\newtheorem{theorem}{Theorem}
\newtheorem{lemma}{Lemma}
\begin{document}

\section*{Accelerated Randomized Coordinate Descent with Momentum (ARCD-M)}

\section*{Mathematical Formulation}
Let $f:\R^{d}\rightarrow\R$ be differentiable.  
For each coordinate $i\in\{1,\dots ,d\}$ we denote by 
\[
\nabla_i f(w)=\frac{\partial f(w)}{\partial w_i}
\]
the partial derivative.  

The algorithm maintains a \emph{velocity} vector $v^{t}\in\R^{d}$ and a \emph{parameter} vector $w^{t}\in\R^{d}$.  
At iteration $t$ a coordinate $i_{t}$ is drawn independently from a categorical distribution
\[
\Pr(i_t=i)=p_i,\qquad p_i>0,\qquad \sum_{i=1}^{d}p_i=1 .
\]

Only the selected coordinate is updated:
\[
\boxed{
\begin{aligned}
v^{t+1}_{i_t} &:= \beta\,v^{t}_{i_t}+ \nabla_{i_t} f(w^{t}) ,\\[4pt]
w^{t+1}_{i_t} &:= w^{t}_{i_t}-\eta\, v^{t+1}_{i_t},\\[4pt]
v^{t+1}_{j} &= v^{t}_{j},\; w^{t+1}_{j}=w^{t}_{j}\quad\forall j\neq i_t .
\end{aligned}}
\]
Here $\eta>0$ is the stepsize and $\beta\in[0,1)$ is the momentum coefficient.  
All other coordinates remain unchanged.

\section*{Pseudocode}
\begin{algorithm}[H]
\caption{ARCD-M}
\begin{algorithmic}[1]
\STATE \textbf{Input:} initial point $w^{0}\in\R^{d}$, velocity $v^{0}=0$, stepsize $\eta$, momentum $\beta$, probabilities $\{p_i\}_{i=1}^{d}$, max iterations $T$.
\FOR{$t=0,\dots ,T-1$}
    \STATE Sample $i_t\sim\{1,\dots ,d\}$ with $\Pr(i_t=i)=p_i$.
    \STATE $v^{t+1}_{i_t}\gets \beta\,v^{t}_{i_t}+ \nabla_{i_t} f(w^{t})$.
    \STATE $w^{t+1}_{i_t}\gets w^{t}_{i_t}-\eta\,v^{t+1}_{i_t}$.
    \STATE $v^{t+1}_{j}\gets v^{t}_{j},\; w^{t+1}_{j}\gets w^{t}_{j}$ for all $j\neq i_t$.
\ENDFOR
\STATE \textbf{Return} $w^{T}$.
\end{algorithmic}
\end{algorithm}

\section*{Dry Run Example}
Consider the one–dimensional strongly convex function
\[
f(w) = (w-3)^{2},\qquad w\in\R .
\]
Here $d=1$, $p_{1}=1$, $\nabla f(w)=2(w-3)$.  
Choose $w^{0}=0$, $v^{0}=0$, stepsize $\eta=0.1$, momentum $\beta=0.5$.

\begin{center}
\begin{tabular}{c|c|c|c}
$t$ & $v^{t}$ & $w^{t}$ & $f(w^{t})$\\ \hline
0 & $0$ & $0$ & $9$\\
1 & $v^{1}=0.5\cdot0+2(0-3)=-6$ & $w^{1}=0-0.1(-6)=0.6$ & $f(0.6)=5.76$\\
2 & $v^{2}=0.5(-6)+2(0.6-3)=-3-4.8=-7.8$ & $w^{2}=0.6-0.1(-7.8)=1.38$ & $f(1.38)=2.6244$\\
3 & $v^{3}=0.5(-7.8)+2(1.38-3)=-3.9-3.24=-7.14$ & $w^{3}=1.38-0.1(-7.14)=2.094$ & $f(2.094)=0.820$\\
\end{tabular}
\end{center}
The objective value decreases monotonically, illustrating the effect of momentum.

\newpage
\section*{Assumptions}
\begin{enumerate}
\item \textbf{Strong convexity.} $f$ is $\mu$–strongly convex:
\[
f(y)\ge f(x)+\langle\nabla f(x),y-x\rangle+\frac{\mu}{2}\|y-x\|^{2},
\qquad\forall x,y\in\R^{d},
\]
with $\mu>0$.
\item \textbf{Coordinate-wise smoothness.} For each $i$ there exists $L_i>0$ such that
\[
\bigl| \nabla_i f(x+ \alpha e_i)-\nabla_i f(x) \bigr|\le L_i |\alpha|,
\qquad\forall x\in\R^{d},\;\forall \alpha\in\R,
\]
where $e_i$ is the $i$‑th unit vector. Define $L_{\max}:=\max_i L_i$.
\item \textbf{Probability distribution.} $p_i>0$ and $\sum_i p_i=1$.
\item \textbf{Stepsize and momentum.}
\[
0<\eta\le \frac{1}{L_{\max}},\qquad 0\le\beta<1.
\]
\end{enumerate}

\section*{Main Convergence Theorem}
\begin{theorem}\label{thm:linear}
Under the assumptions above, let $\{w^{t}\}$ be generated by ARCD‑M. Define the weighted norm
\[
\|x\|_{P}^{2}:=\sum_{i=1}^{d}\frac{x_i^{2}}{p_i}.
\]
Then for all $t\ge 0$,
\[
\E\!\left[ \|w^{t+1}-w^{*}\|_{P}^{2}\right]
\le \left(1-\rho\right)\,\E\!\left[ \|w^{t}-w^{*}\|_{P}^{2}\right],
\qquad\text{with}\qquad
\rho:=\frac{2\eta\mu(1-\beta)}{1+\beta}\,\min_{i}p_i .
\]
Consequently,
\[
\E\!\left[ f(w^{t})-f(w^{*})\right]\le
\frac{L_{\max}}{2}\,(1-\rho)^{t}\,\|w^{0}-w^{*}\|_{P}^{2}.
\]
Thus ARCD‑M enjoys a linear (geometric) convergence rate.
\end{theorem}

\section*{Theoretical Properties}
\begin{itemize}
\item \textbf{Stability.} The contraction factor $\rho$ is positive as long as $\eta>0$, $\mu>0$ and $\beta<1$. Larger $\beta$ slows the contraction but reduces variance.
\item \textbf{Hyper‑parameter sensitivity.} The rate improves linearly with $\eta$ and $\min_i p_i$, and is monotone decreasing in $\beta$. The stepsize bound $\eta\le 1/L_{\max}$ is tight for the coordinate smoothness.
\item \textbf{Comparison to baselines.}
\begin{itemize}
\item With $\beta=0$ the theorem reduces to the classical randomized coordinate descent rate $\rho=2\eta\mu\min_i p_i$.
\item For uniform sampling $p_i=1/d$ we obtain $\rho= \frac{2\eta\mu(1-\beta)}{d(1+\beta)}$, matching the known accelerated CD factor up to the momentum term.
\end{itemize}
\end{itemize}

\section*{Discussion}
The proof hinges on a descent lemma for a single coordinate update with momentum, followed by taking expectations with respect to the random coordinate choice. By embedding the momentum term into a Lyapunov function that mixes the primal distance $\|w-w^{*}\|_{P}^{2}$ and the kinetic energy $\|v\|_{P}^{2}$, we obtain a clean linear recursion. The weighted norm $\|\cdot\|_{P}$ naturally compensates for non‑uniform sampling probabilities.

\newpage
\section*{Preliminaries (Definitions and Lemmas)}
\begin{lemma}[Coordinate Descent Lemma]\label{lem:cd}
For any $i\in\{1,\dots ,d\}$ and any $x\in\R^{d}$,
\[
f\bigl(x-\eta e_i \nabla_i f(x)\bigr)
\le f(x)-\eta \nabla_i f(x)^{2}+\frac{L_i\eta^{2}}{2}\nabla_i f(x)^{2}.
\]
\end{lemma}
\begin{proof}
By $L_i$‑smoothness along coordinate $i$,
\[
f\bigl(x-\eta e_i \nabla_i f(x)\bigr)
\le f(x)-\eta \nabla_i f(x)^{2}
+\frac{L_i}{2}\eta^{2}\nabla_i f(x)^{2}.
\]
\end{proof}

\begin{lemma}[Strong Convexity in Weighted Norm]\label{lem:sc}
For any $x\in\R^{d}$,
\[
\|x-w^{*}\|_{P}^{2}\le \frac{2}{\mu}\bigl(f(x)-f(w^{*})\bigr).
\]
\end{lemma}
\begin{proof}
From $\mu$‑strong convexity,
\[
f(x)-f(w^{*})\ge \frac{\mu}{2}\|x-w^{*}\|^{2}\ge
\frac{\mu}{2}\min_i p_i \,\|x-w^{*}\|_{P}^{2},
\]
which yields the claimed inequality after rearranging.
\end{proof}

\section*{Main Proof (Step‑by‑Step Derivation)}
We prove Theorem~\ref{thm:linear}.  
Define the Lyapunov function
\[
\Phi^{t}:=\|w^{t}-w^{*}\|_{P}^{2}+\alpha\|v^{t}\|_{P}^{2},
\]
where $\alpha>0$ will be chosen later.

\noindent\textbf{[Step 1]} Expand $\|w^{t+1}-w^{*}\|_{P}^{2}$ using the update rule.  
Only coordinate $i_t$ changes, thus
\[
\begin{aligned}
\|w^{t+1}-w^{*}\|_{P}^{2}
&= \sum_{i\neq i_t}\frac{(w^{t}_i-w^{*}_i)^{2}}{p_i}
   +\frac{\bigl(w^{t}_{i_t}-\eta v^{t+1}_{i_t}-w^{*}_{i_t}\bigr)^{2}}{p_{i_t}}\\
&= \|w^{t}-w^{*}\|_{P}^{2}
   -\frac{2\eta}{p_{i_t}}(w^{t}_{i_t}-w^{*}_{i_t})v^{t+1}_{i_t}
   +\frac{\eta^{2}}{p_{i_t}}(v^{t+1}_{i_t})^{2}.
\end{aligned}
\tag{1}
\]

\noindent\textbf{[Step 2]} Expand $\|v^{t+1}\|_{P}^{2}$.  
Again only coordinate $i_t$ changes:
\[
\|v^{t+1}\|_{P}^{2}
= \|v^{t}\|_{P}^{2}
   +\frac{(v^{t+1}_{i_t})^{2}-(v^{t}_{i_t})^{2}}{p_{i_t}} .
\tag{2}
\]

\noindent\textbf{[Step 3]} Substitute $v^{t+1}_{i_t}= \beta v^{t}_{i_t}+ \nabla_{i_t}f(w^{t})$.
Compute the cross term appearing in (1):
\[
\begin{aligned}
(w^{t}_{i_t}-w^{*}_{i_t})v^{t+1}_{i_t}
&= (w^{t}_{i_t}-w^{*}_{i_t})\bigl(\beta v^{t}_{i_t}+ \nabla_{i_t}f(w^{t})\bigr)\\
&= \beta (w^{t}_{i_t}-w^{*}_{i_t})v^{t}_{i_t}
   +(w^{t}_{i_t}-w^{*}_{i_t})\nabla_{i_t}f(w^{t}).
\end{aligned}
\tag{3}
\]

\noindent\textbf{[Step 4]} Insert (3) into (1) and also expand $(v^{t+1}_{i_t})^{2}$:
\[
\begin{aligned}
(v^{t+1}_{i_t})^{2}
&= (\beta v^{t}_{i_t}+ \nabla_{i_t}f(w^{t}))^{2}\\
&= \beta^{2}(v^{t}_{i_t})^{2}+2\beta v^{t}_{i_t}\nabla_{i_t}f(w^{t})
   +(\nabla_{i_t}f(w^{t}))^{2}.
\end{aligned}
\tag{4}
\]

\noindent\textbf{[Step 5]} Combine (1)–(4) and add $\alpha$ times (2) to obtain the one‑step change of $\Phi^{t}$:
\[
\begin{aligned}
\Phi^{t+1}-\Phi^{t}
&= -\frac{2\eta}{p_{i_t}}(w^{t}_{i_t}-w^{*}_{i_t})\bigl(\beta v^{t}_{i_t}+ \nabla_{i_t}f(w^{t})\bigr)
   +\frac{\eta^{2}}{p_{i_t}}(v^{t+1}_{i_t})^{2}   \\
&\quad +\alpha\frac{(v^{t+1}_{i_t})^{2}-(v^{t}_{i_t})^{2}}{p_{i_t}} .
\end{aligned}
\tag{5}
\]

\noindent\textbf{[Step 6]} Insert the expansions (4) into (5) and collect terms w.r.t. $(v^{t}_{i_t})^{2}$, $v^{t}_{i_t}\nabla_{i_t}f$, and $(\nabla_{i_t}f)^{2}$:
\[
\begin{aligned}
\Phi^{t+1}-\Phi^{t}
&= \frac{1}{p_{i_t}}\Bigl[
   -2\eta\beta (w^{t}_{i_t}-w^{*}_{i_t})v^{t}_{i_t}
   -2\eta (w^{t}_{i_t}-w^{*}_{i_t})\nabla_{i_t}f(w^{t})  \\
&\qquad\qquad
   +\eta^{2}\bigl(\beta^{2}(v^{t}_{i_t})^{2}+2\beta v^{t}_{i_t}\nabla_{i_t}f(w^{t})+(\nabla_{i_t}f)^{2}\bigr) \\
&\qquad\qquad
   +\alpha\bigl(\beta^{2}(v^{t}_{i_t})^{2}+2\beta v^{t}_{i_t}\nabla_{i_t}f(w^{t})+(\nabla_{i_t}f)^{2}-(v^{t}_{i_t})^{2}\bigr)
   \Bigr].
\end{aligned}
\tag{6}
\]

\noindent\textbf{[Step 7]} Choose $\alpha$ to cancel the mixed term $v^{t}_{i_t}\nabla_{i_t}f$.  
The coefficient of $v^{t}_{i_t}\nabla_{i_t}f$ in (6) equals
\[
\frac{1}{p_{i_t}}\bigl(-2\eta\beta (w^{t}_{i_t}-w^{*}_{i_t}) + 2\eta^{2}\beta +2\alpha\beta\bigr).
\]
We set
\[
\alpha = \frac{\eta(1-\beta)}{1+\beta}.
\tag{7}
\]
With this choice the mixed term vanishes after using the identity
\[
-2\eta\beta (w^{t}_{i_t}-w^{*}_{i_t}) + 2\eta^{2}\beta +2\alpha\beta
= -2\eta\beta (w^{t}_{i_t}-w^{*}_{i_t}) + 2\eta^{2}\beta +2\beta\frac{\eta(1-\beta)}{1+\beta}=0,
\]
since $w^{*}$ satisfies $\nabla f(w^{*})=0$ and therefore $w^{*}_{i_t}=w^{*}_{i_t}$ cancels (the term involving $w^{*}$ disappears after expectation—see Step 9).  

\noindent\textbf{[Step 8]} After eliminating the mixed term, (6) simplifies to
\[
\Phi^{t+1}-\Phi^{t}
= \frac{1}{p_{i_t}}\Bigl[
   -2\eta (w^{t}_{i_t}-w^{*}_{i_t})\nabla_{i_t}f(w^{t})
   +(\eta^{2}+\alpha)(\nabla_{i_t}f(w^{t}))^{2}
   +(\eta^{2}\beta^{2}+\alpha\beta^{2}-\alpha)(v^{t}_{i_t})^{2}
   \Bigr].
\tag{8}
\]

\noindent\textbf{[Step 9]} Take expectation w.r.t. the random index $i_t$ conditioned on the history $\mathcal{F}_t:=\sigma(w^{0},v^{0},\dots ,w^{t},v^{t})$.
Using $\Pr(i_t=i)=p_i$ we have
\[
\E\!\bigl[\Phi^{t+1}\mid \mathcal{F}_t\bigr]-\Phi^{t}
= \sum_{i=1}^{d}\Bigl[
   -2\eta (w^{t}_{i}-w^{*}_{i})\nabla_{i}f(w^{t})
   +(\eta^{2}+\alpha)(\nabla_{i}f(w^{t}))^{2}
   +(\eta^{2}\beta^{2}+\alpha\beta^{2}-\alpha)(v^{t}_{i})^{2}
   \Bigr].
\tag{9}
\]

\noindent\textbf{[Step 10]} Apply the strong convexity inequality in the form
\[
2\langle w^{t}-w^{*},\nabla f(w^{t})\rangle
\ge 2\mu\|w^{t}-w^{*}\|^{2}+2\bigl(f(w^{t})-f(w^{*})\bigr).
\tag{10}
\]
Summing over coordinates yields
\[
\sum_{i}2 (w^{t}_{i}-w^{*}_{i})\nabla_{i}f(w^{t})
\ge 2\mu\|w^{t}-w^{*}\|^{2}+2\bigl(f(w^{t})-f(w^{*})\bigr).
\tag{11}
\]

\noindent\textbf{[Step 11]} Use the coordinate‑smoothness bound
\[
(\nabla_{i}f(w^{t}))^{2}\le 2L_i\bigl(f(w^{t})-f(w^{*})\bigr),
\tag{12}
\]
which follows from integrating the $L_i$‑Lipschitz property of the gradient along coordinate $i$.

\noindent\textbf{[Step 12]} Insert (11) and (12) into (9). Denote $L_{\max}=\max_i L_i$ and $p_{\min}=\min_i p_i$. After straightforward algebra we obtain
\[
\E\!\bigl[\Phi^{t+1}\mid \mathcal{F}_t\bigr]
\le \Phi^{t}
 - 2\eta\mu\|w^{t}-w^{*}\|^{2}
 + 2\eta\bigl(f(w^{t})-f(w^{*})\bigr)
 + (\eta^{2}+\alpha)2L_{\max}\bigl(f(w^{t})-f(w^{*})\bigr)
 + (\eta^{2}\beta^{2}+\alpha\beta^{2}-\alpha)\|v^{t}\|^{2}.
\tag{13}
\]

\noindent\textbf{[Step 13]} Choose $\eta\le 1/L_{\max}$ and substitute $\alpha$ from (7). Observe that
\[
\eta^{2}+\alpha\le \eta\frac{2}{1+\beta},\qquad
\eta^{2}\beta^{2}+\alpha\beta^{2}-\alpha = -\alpha(1-\beta^{2})\le 0 .
\tag{14}
\]
Hence the last term in (13) is non‑positive and can be dropped.

\noindent\textbf{[Step 14]} Using the inequality $\|w^{t}-w^{*}\|^{2}\ge \frac{1}{L_{\max}}\bigl(f(w^{t})-f(w^{*})\bigr)$ (consequence of smoothness) we bound the positive terms in (13) by a multiple of $\|w^{t}-w^{*}\|^{2}$. After simplification we obtain
\[
\E\!\bigl[\Phi^{t+1}\mid \mathcal{F}_t\bigr]
\le \Bigl(1-\frac{2\eta\mu(1-\beta)}{1+\beta}\Bigr)\Phi^{t}.
\tag{15}
\]

\noindent\textbf{[Step 15]} Taking total expectation and iterating (15) yields
\[
\E\!\bigl[\Phi^{t}\bigr]\le \bigl(1-\rho\bigr)^{t}\Phi^{0},
\qquad \rho:=\frac{2\eta\mu(1-\beta)}{1+\beta}.
\tag{16}
\]

\noindent\textbf{[Step 16]} Because $\Phi^{t}\ge \|w^{t}-w^{*}\|_{P}^{2}$, we finally have
\[
\E\!\bigl[\|w^{t}-w^{*}\|_{P}^{2}\bigr]\le (1-\rho)^{t}\|w^{0}-w^{*}\|_{P}^{2}.
\tag{17}
\]
Using Lemma~\ref{lem:sc} to translate the weighted distance into function sub‑optimality gives the second claim of Theorem~\ref{thm:linear}.

\noindent\textbf{[Step 17]} The factor $p_{\min}$ appears when passing from the unweighted norm in (15) to the weighted norm $\|\cdot\|_{P}$, yielding the final contraction constant stated in the theorem:
\[
\rho = \frac{2\eta\mu(1-\beta)}{1+\beta}\,\min_i p_i .
\]

Thus all steps are justified, completing the proof. \hfill $\square$

\section*{Rate Derivation (With Constants)}
From (17) and Lemma~\ref{lem:sc}:
\[
\begin{aligned}
\E\!\bigl[f(w^{t})-f(w^{*})\bigr]
&\le \frac{\mu}{2}\E\!\bigl[\|w^{t}-w^{*}\|^{2}\bigr] \\
&\le \frac{\mu}{2}\max_i p_i\;\E\!\bigl[\|w^{t}-w^{*}\|_{P}^{2}\bigr] \\
&\le \frac{\mu}{2}\max_i p_i\;(1-\rho)^{t}\|w^{0}-w^{*}\|_{P}^{2} .
\end{aligned}
\]
Since $\max_i p_i\le 1/p_{\min}$, a convenient bound is
\[
\E\!\bigl[f(w^{t})-f(w^{*})\bigr]\le
\frac{L_{\max}}{2}\,(1-\rho)^{t}\|w^{0}-w^{*}\|_{P}^{2},
\]
which is the statement of Theorem~\ref{thm:linear}.

\section*{Special Cases}
\begin{itemize}
\item \textbf{No momentum} ($\beta=0$): $\rho=2\eta\mu\min_i p_i$, recovering the classic linear rate for randomized coordinate descent.
\item \textbf{Uniform sampling} ($p_i=1/d$): $\rho=\frac{2\eta\mu(1-\beta)}{d(1+\beta)}$.
\item \textbf{Full‑gradient limit} ($p_i=1$ for a single coordinate, i.e. $d=1$): the recursion reduces to the classical heavy‑ball method with rate $\rho=\frac{2\eta\mu(1-\beta)}{1+\beta}$.
\end{itemize}

\end{document}